\clearpage\SourceLeaf{206}
\selectlanguage{ngerman}\fancyhead[L]{\small ANZEIGE.}\pdfbookmark[0]{ANZEIGEN.}{anzeigen206}
\begin{center}\textbf{A N Z E I G E N.}\end{center}
\par\Needspace{9\baselineskip}\begin{center}
\rule{0.52\linewidth}{0.3pt}\par\medskip
Göttingische gelehrte Anzeigen.\quad 1812 Februar 10.\par\medskip
\rule{0.52\linewidth}{0.3pt}
\end{center}

Die logarithmischen und Kreisfunktionen sind als die einfachsten Arten der transzendenten Funktionen diejenigen, mit denen sich die Analytiker am meisten beschäftigt haben. Sie verdienten diese Ehre sowohl wegen ihres ständigen Eingreifens in fast alle theoretischen und praktischen mathematischen Untersuchungen als auch wegen des nahezu unerschöpflichen Reichtums an interessanten Wahrheiten, den ihre Theorie bietet.\quad Weit weniger wurden bisher andere transzendente Funktionen untersucht, die sich nicht auf jene zurückführen lassen, sondern als eigene höhere Arten betrachtet werden müssen, die nur in besonderen Fällen mit jenen zusammenhängen. Dennoch sind manche dieser Funktionen nicht weniger reich an interessanten Beziehungen und daher für denjenigen, der die Analysis um ihrer selbst willen schätzt, nicht weniger wichtig; ebenso muss ihr häufiges Vorkommen in verschiedenen anderen Untersuchungen sie demjenigen empfehlen, der gern zuerst nach dem praktischen Nutzen fragt. Professor \textsc{Gauss} beschäftigt sich bereits seit vielen Jahren mit Untersuchungen über solche höheren transzendenten Funktionen, deren weitreichende Ergebnisse das Gesagte bestätigen. Einen im Verhältnis dazu freilich nur sehr kleinen Teil dieser Untersuchungen, der gleichsam als Einleitung zu einer künftig vorzulegenden Reihe von Abhandlungen angesehen werden kann, hat er am 30. Januar unter dem Titel

\begin{center}\emph{Disquisitiones generales circa seriem infinitam}\end{center}
\[
\begin{aligned}&1+\frac{\alpha.\beta}{1.\gamma}x
+\frac{\alpha.\alpha+1.\beta.\beta+1}{1.2.\gamma.\gamma+1}xx
\\[5pt]&+\frac{\alpha.\alpha+1.\alpha+2.\beta.\beta+1.\beta+2}{1.2.3.\gamma.\gamma+1.\gamma+2}x^3+\mathit{usw.}\end{aligned}
\]
\begin{center}\emph{Pars prior,}\end{center}
\SourceLeaf{207}
der königlichen Gesellschaft der Wissenschaften als Vorlesung übergeben; die Hauptpunkte ihres Inhalts wollen wir hier anzeigen.

Die transzendenten Funktionen haben ihre wahre Quelle stets, offen oder verborgen, im Unendlichen. Die Operationen des Integrierens, der Summation unendlicher Reihen, der Entwicklung unendlicher Produkte, der ins Unendliche fortlaufenden Kettenbrüche oder überhaupt die Annäherung an eine Grenze durch Operationen, die nach bestimmten Gesetzen ohne Ende fortgesetzt werden --- dies ist der eigentliche Boden, auf dem die transzendenten Funktionen erzeugt werden; oder, wenn man lieber ein anderes Bild verwenden möchte, dies sind die eigentlichen Wege, auf denen man zu ihnen gelangt. Zu \emph{einem} Ziel führen gewöhnlich mehrere solcher Wege: Die Umstände und die Ziele, die man sich setzt, müssen bestimmen, welchen man zuerst oder vorzugsweise wählen will. Die Reihe, die den Gegenstand der vorliegenden Abhandlung bildet, besitzt eine sehr umfassende Allgemeinheit. Je nachdem, wie die Größen $\alpha$, $\beta$, $\gamma$ bestimmt werden (die der Verfasser zusammen mit $x$ durch die Bezeichnungen erstes, zweites, drittes und viertes \emph{Element} unterscheidet, während er der Kürze halber die gesamte durch die Reihe dargestellte Funktion mit den Zeichen $F(\alpha,\beta,\gamma,x)$ bezeichnet), stellt sie algebraische, logarithmische, trigonometrische oder höhere transzendente Funktionen dar. Man kann behaupten, dass bisher kaum eine transzendente Funktion von den Analytikern untersucht wurde, die sich nicht auf diese Reihe zurückführen ließe. Eine große Anzahl bereits gefundener Wahrheiten über solche schon betrachteten Funktionen lässt sich aus der allgemeinen Natur der durch unsere Reihen dargestellten Funktion ableiten. Schon deshalb würden Untersuchungen darüber die Aufmerksamkeit der Geometer verdienen, obwohl dies nur als Nebensache und die Eröffnung des Zugangs zu neuen Wahrheiten als Hauptzweck anzusehen ist. Die vorliegende Abhandlung enthält erst die Hälfte der allgemeinen Untersuchungen des Verfassers, deren zu großer Umfang eine Teilung notwendig machte. Hier gilt die Reihe selbst als Ursprung der transzendenten Funktionen, die im weiteren Verlauf der Arbeit aus einer allgemeiner anwendbaren Quelle abgeleitet und von einem höheren Standpunkt aus betrachtet werden sollen. Die erste Erzeugungsart erfordert ihrer Natur nach die Beschränkung auf die Fälle, in denen die Reihe konvergiert, also auf diejenigen, in denen das vierte Element $x$, positiv oder negativ,
\SourceLeaf{208}
den Wert $1$ nicht überschreitet: Die andere Erzeugungsart wird diese Beschränkung beseitigen. Doch bereits die erste Erzeugungsart führt auf einem bequemeren und gleichsam elementareren Weg zu einer Vielzahl bemerkenswerter Wahrheiten; deshalb hat der Verfasser mit ihr begonnen.

Diese erste Hälfte der Untersuchungen gliedert sich in drei Abschnitte, denen einige allgemeine Bemerkungen vorangestellt sind. Um eine Probe der weitreichenden Anwendbarkeit der Reihe zu geben, werden zunächst $23$ verschiedene Reihenentwicklungen algebraischer, logarithmischer und trigonometrischer Funktionen auf sie zurückgeführt. Dasselbe geschieht als Beispiel höherer transzendenter Funktionen mit den Koeffizienten der Reihe, die aus der Entwicklung von $(aa+bb-2ab\cos\varphi)^n$ entsteht und nach den Kosinus der Vielfachen von $\varphi$ fortschreitet, und zwar bei diesen Koeffizienten auf drei verschiedene Arten.

Der \textit{erste Abschnitt} befasst sich mit den Beziehungen zwischen Reihen der obigen Form, bei denen sich die Werte eines der ersten drei Elemente um eine Einheit unterscheiden, während die Werte der drei übrigen gleich sind. Solche Reihen nennt der Verfasser \textit{series contiguae}; im Deutschen könnte man sie etwa \textit{verwandte Reihen} nennen. Jeder Reihe $F(\alpha,\beta,\gamma,x)$ stehen somit sechs verwandte zur Seite, nämlich $F(\alpha+1,\beta,\gamma,x)$, $F(\alpha-1,\beta,\gamma,x)$, $F(\alpha,\beta+1,\gamma,x)$, $F(\alpha,\beta-1,\gamma,x)$, $F(\alpha,\beta,\gamma+1,x)$, $F(\alpha,\beta,\gamma-1,x)$. Hier wird gezeigt, dass zwischen der ersten und je zweien der verwandten eine lineare Gleichung besteht. Auf diese Weise entstehen fünfzehn Gleichungen. Daraus folgt der wichtige Satz, dass, wenn $\alpha'-\alpha$, $\alpha''-\alpha$, $\beta'-\beta$, $\beta''-\beta$, $\gamma'-\gamma$, $\gamma''-\gamma$ ganze Zahlen sind, auch zwischen $F(\alpha,\beta,\gamma,x)$, $F(\alpha',\beta',\gamma',x)$, $F(\alpha'',\beta'',\gamma'',x)$ eine lineare Gleichung besteht und somit, allgemein gesprochen, aus den Werten zweier dieser Funktionen der Wert der dritten abgeleitet werden kann. Einige der einfachsten oder anderweitig bemerkenswerten Fälle hat der Verfasser hier außerdem gesondert zusammengestellt.

Der \textit{zweite Abschnitt} liefert Umwandlungen in Kettenbrüche, und zwar für die Quotienten
\[
\frac{F(\alpha,\beta+1,\gamma+1,x)}{F(\alpha,\beta,\gamma,x)},\quad
\frac{F(\alpha,\beta+1,\gamma,x)}{F(\alpha,\beta,\gamma,x)},\quad
\frac{F(\alpha-1,\beta+1,\gamma,x)}{F(\alpha,\beta,\gamma,x)}
\]
auf die sich durch die offenbar zulässige Vertauschung der beiden ersten Elemente noch drei andere zurückführen lassen. Fast alle bisher bekannten Entwicklungen in Kettenbrüche sind nur besondere Fälle dieser Sätze.
\SourceLeaf{209}
Besonders bemerkenswert ist der Fall, in dem man in der zweiten Entwicklung $\beta=0$ setzt. Daraus folgt ein Satz, den wir wegen seiner weitreichenden Anwendbarkeit hier beifügen. Die Funktion $F(\alpha,1,\gamma.x)$ oder, was dasselbe ist, die Reihe
\[
1+\frac{\alpha}{\gamma}x
+\frac{\alpha\,.\,\alpha+1}{\gamma\,.\,\gamma+1}xx
+\frac{\alpha\,.\,\alpha+1\,.\,\alpha+2}{\gamma\,.\,\gamma+1\,.\,\gamma+2}x^3+\mathrm{usw.}
\]
ergibt den Kettenbruch
\[
\dfrac{1}{\displaystyle 1-\dfrac{ax}{\displaystyle 1-\dfrac{bx}{\displaystyle 1-\dfrac{cx}{\displaystyle 1-\dfrac{dx}{1-\mathrm{usw.}}}}}}
\]
in dem die Koeffizienten $a,b,c,d$ usw. nach folgendem Gesetz fortschreiten:
\[
\begin{array}{ll}
a=\dfrac{\alpha}{\gamma}&b=\dfrac{\gamma-\alpha}{\gamma(\gamma+1)}\\[6pt]
c=\dfrac{(\alpha+1)\gamma}{(\gamma+1)(\gamma+2)}&d=\dfrac{2(\gamma+1-\alpha)}{(\gamma+2)(\gamma+3)}\\[6pt]
e=\dfrac{(\alpha+2)(\gamma+1)}{(\gamma+3)(\gamma+4)}&f=\dfrac{3(\gamma+2-\alpha)}{(\gamma+4)(\gamma+5)}\quad\mbox{usw.}
\end{array}
\]
Hiernach lassen sich zum Beispiel die Potenz eines Binoms, die Reihen für $\log(1+x)$, $\log\dfrac{1+x}{1-x}$, für Exponentialgrößen, für den Bogen aus der Tangente oder aus dem Sinus und andere in unendliche Kettenbrüche umwandeln. Darauf beruhen auch die in der \textit{Theoria motus corporum coelestium} angegebenen Umwandlungen in solche Brüche, deren Beweise der Verfasser hier nachliefert.

Den weitaus größten Teil der Abhandlung nimmt der \textit{dritte Abschnitt} ein, der den Wert der Reihe behandelt, wenn man das vierte Element $=1$ setzt. Nachdem zunächst mit geometrischer Strenge bewiesen wurde, dass die Reihe für $x=1$ nur dann zu einer endlichen Summe konvergiert, wenn $\gamma-\alpha-\beta$ eine positive Größe ist, führt der Verfasser diese Summe oder $F(\alpha,\beta,\gamma,1)$ auf den Ausdruck $\dfrac{\Pi(\gamma-1)\,.\,\Pi(\gamma-\alpha-\beta-1)}{\Pi(\gamma-\alpha-1)\,.\,\Pi(\gamma-\beta-1)}$ zurück. Dabei bezeichnet das Symbol $\Pi$ eine eigene Art transzendenter Funktionen, deren Erzeugung der Verfasser auf ein unendliches Produkt gründet. Diese in der gesamten Analysis höchst wichtige Funktion ist im Grunde nichts anderes als \textsc{Eulers} inexplikable Funktion
\[
\Pi z=1\,.\,2\,.\,3\,.\,4\,.\,.\,.\,.\,.\,z
\]
\SourceLeaf{210}
Doch \textit{diese} Erzeugungsart oder Definition ist nach dem Urteil des Verfassers völlig unstatthaft, da sie nur für positive ganzzahlige Werte von $z$ einen klaren Sinn hat. Die vom Verfasser gewählte Begründungsart ist allgemein anwendbar und ergibt selbst bei imaginären Werten von $z$ einen ebenso klaren Sinn wie bei reellen. Dabei besteht keinerlei Gefahr, in solche Paradoxien und Widersprüche zu geraten wie ehemals Herr\ \textsc{Kramp} bei seinen numerischen Fakultäten. Diese lassen sich, wie man leicht zeigen kann, auf die obige Funktion zurückführen, scheinen aber für die Aufnahme in die Analysis weniger geeignet als diese, da sie von drei Größen abhängen, während diese nur von einer abhängt und dennoch als ebenso allgemein angesehen werden muss. Der Verfasser wünscht, dass dieser transzendenten Funktion $\Pi z$ in der Analysis das Bürgerrecht gegeben wird. Vielleicht würde die Wahl eines eigenen Namens für sie dies am meisten fördern: Das Recht dazu mag demjenigen vorbehalten bleiben, der die wichtigsten Entdeckungen in der Theorie dieser Funktion macht, die der Anstrengungen der Geometer sehr würdig ist. Der Verfasser hat hier bereits eine beträchtliche Anzahl bemerkenswerter Sätze über sie zusammengestellt, von denen ein Teil als neu anzusehen ist. Der Raum erlaubt uns nicht, hier auf ihre Einzelheiten einzugehen: Wir heben nur den einen hervor, dass sich der Wert des Integrals $\int x^{\lambda-1}(1-x^\mu)^\nu\,\mathrm{d}x$ von $x=0$ bis $x=1$ leicht auf die Funktion $\Pi$ zurückführen lässt und dass sich alle von \textsc{Euler} für solche Integrale zum Teil mühsam gefundenen Beziehungen mit größter Leichtigkeit aus den allgemeinen Eigenschaften jener Funktionen ableiten lassen. Umgekehrt lässt sich stets $\Pi z$, wenn $z$ eine rationale Größe ist, durch einige solche bestimmte Integrale darstellen.

Nicht weniger bemerkenswert ist die aus der Differentiation von $\Pi z$ entstehende, ebenfalls transzendente Funktion oder vielmehr
\[
\frac{\mathrm{d}\log\Pi z}{\mathrm{d}z}=\frac{\mathrm{d}\Pi z}{\Pi z\,.\,\mathrm{d}z}
\]
die der Verfasser mit $\Psi z$ bezeichnet hat und die ebenfalls eine besondere \textit{Benennung} verdiente. Von den zahlreichen bemerkenswerten Eigenschaften dieser Funktion, die in der Abhandlung aufgestellt sind, führen wir hier nur die eine an: Allgemein lässt sich $\Psi z-\Psi 0$, wenn $z$ eine rationale Größe ist, auf Logarithmen und Kreisfunktionen zurückführen. $\Psi 0$ selbst ist aber die bekannte, von \textsc{Euler} und anderen untersuchte Zahl $0{,}5772156649\,.\,.\,.\,.$, negativ genommen. Der Verfasser teilt sie hier nach einer von ihm selbst durchgeführten Rechnung \SourceLeaf{211}%
auf 23 Dezimalstellen mit, deren letzte etwas von \textsc{Mascheroni}s Bestimmung abweichen. Im Übrigen hängen sowohl $\Pi z$ als auch $\Psi z$ mit mehreren bemerkenswerten Integralen für bestimmte Werte der veränderlichen Größe zusammen.

Diesem dritten Abschnitt ist außerdem eine Tafel für $\log\Pi z$ und für $\Psi z$ beigefügt, die unter der Aufsicht von Professor \textsc{Gauss} von Herrn\ \textsc{Nicolai} mit größter Sorgfalt berechnet wurde und in der das Argument $z$ alle einzelnen Hundertstel von 0 bis 1 durchläuft. Aus der Theorie dieser Funktionen ist klar, dass sich alle anderen Werte leicht auf diese Werte von $z$ zurückführen lassen.

